% GENERATED FILE. Edit canonical lessons, then run npm run generate:books.
% canonical-source-hash: fnv1a64:19d69d9a8c33f239
% canonical-lessons: LA-C445-deverticulum, LA-C445-deversor, LA-C445-pernocto, LA-C445-stabularius, LA-C445-hospitalis

\chapter{Roadside inn, To lodge, To spend the night, Stable-keeper, Hospitable}
\label{ch:la-roadside-inn-to-lodge-to-spend-the-night-stable-keeper-hospitable}
\hlchaptermodality{445}

\begin{chapteropening}
\textbf{By the end of this chapter:} \emph{I can talk about stopping at an inn: a roadside inn, lodging there, spending the night, the stable-keeper and a hospitable host.}

Complete the last lesson of chapter 445: I can talk about stopping at an inn: a roadside inn, lodging there, spending the night, the stable-keeper and a hospitable host.
\end{chapteropening}

\section[dēverticulum, dēverticulī]{dēverticulum, dēverticulī — a roadside inn, a lodging; a side road}

\label{lesson:LA-C445-deverticulum}



\begin{quote}
Before the new one: say the Latin for of the shore, then the Latin for to sail out.
\end{quote}

\subsection*{You'll want to know: dēverticulum, dēverticulī}
\textbf{dēverticulum, dēverticulī} — \textquotedblleft{}a roadside inn, a lodging; a side road\textquotedblright{}.

\begin{quote}
Nox venit, itaque viātōrēs in dēverticulum intrant.
\end{quote}

\begin{quote}
\emph{Night is coming, so the travellers go into a roadside inn.}
\end{quote}

Second declension, neuter: \emph{dēverticulum, dēverticulī n.} A place you turn off the road for: first a side road, then the inn at its end.

\begin{cousinweb}
\emph{dēverticulum} is built on \emph{dēvertō} 'to turn aside' with the suffix \emph{-culum} 'place for' (EtymDB; Lewis \& Short): a place you turn aside to. \emph{dēvertō} comes from \emph{vertō} 'to turn'.
\end{cousinweb}

\begin{grammarlens}[title={What you've built}]
One more word for telling, explaining and describing at B1.
\end{grammarlens}

\subsection*{Your turn}
\begin{itemize}
\raggedright
  \item \emph{Say it:} \emph{dēverticulum, dēverticulī}
  \item \emph{Say it:} \emph{dēverticulum, dēverticulī}, once more
  \item \emph{Say it:} a sentence of your own with \emph{dēverticulum, dēverticulī}
  \item \emph{Recall:} say the Latin for of the shore, then the Latin for to sail out, then say \emph{dēverticulum, dēverticulī} again
\end{itemize}

\begin{tcolorbox}[breakable,colback=teal!4,colframe=teal!35!black,title={Before you move on}]
What does dēverticulum, dēverticulī mean? (\textquotedblleft{}A roadside inn, a lodging; a side road\textquotedblright{}.) Say it once more.
\end{tcolorbox}
\section[dēversor, dēversārī]{dēversor, dēversārī — to lodge, to stay as a guest}

\label{lesson:LA-C445-deversor}



\begin{quote}
Before the new one: say the Latin for to sail out, then the Latin for a roadside inn.
\end{quote}

\subsection*{You'll want to know: dēversor, dēversārī}
\textbf{dēversor, dēversārī} — \textquotedblleft{}to lodge, to stay as a guest\textquotedblright{}.

\begin{quote}
Mercātor apud amīcum in oppidō dēversātur.
\end{quote}

\begin{quote}
\emph{The merchant lodges with a friend in the town.}
\end{quote}

Deponent: \emph{dēversor, dēversārī, dēversātus sum}, often with \emph{apud} 'at the house of'. The inn is a \emph{dēversōrium}.

\begin{cousinweb}
\emph{dēversor} is a frequentative built on \emph{dēvertō} 'to turn aside' (Lewis \& Short): to keep turning in, to stay as a guest. \emph{dēversōrium} 'inn' comes from the same verb.
\end{cousinweb}

\begin{grammarlens}[title={What you've built}]
One more word for telling, explaining and describing at B1.
\end{grammarlens}

\subsection*{Your turn}
\begin{itemize}
\raggedright
  \item \emph{Say it:} \emph{dēversor, dēversārī}
  \item \emph{Say it:} \emph{dēversor, dēversārī}, once more
  \item \emph{Say it:} a sentence of your own with \emph{dēversor, dēversārī}
  \item \emph{Recall:} say the Latin for to sail out, then the Latin for a roadside inn, then say \emph{dēversor, dēversārī} again
\end{itemize}

\begin{tcolorbox}[breakable,colback=teal!4,colframe=teal!35!black,title={Before you move on}]
What does dēversor, dēversārī mean? (\textquotedblleft{}To lodge, to stay as a guest\textquotedblright{}.) Say it once more.
\end{tcolorbox}
\section[pernoctō, pernoctāre]{pernoctō, pernoctāre — to spend the night, to stay all night}

\label{lesson:LA-C445-pernocto}



\begin{quote}
Before the new one: say the Latin for a roadside inn, then the Latin for to lodge.
\end{quote}

\subsection*{You'll want to know: pernoctō, pernoctāre}
\textbf{pernoctō, pernoctāre} — \textquotedblleft{}to spend the night, to stay all night\textquotedblright{}.

\begin{quote}
Viātōrēs in stabulō pernoctant, quod caupōna plēna est.
\end{quote}

\begin{quote}
\emph{The travellers spend the night in the stable, because the inn is full.}
\end{quote}

\emph{pernoctō, pernoctāre, pernoctāvī, pernoctātum}. From \emph{per} 'through' and \emph{nox} 'night'.

\begin{cousinweb}
\emph{pernoctō} is built from \emph{per} 'through' and \emph{nox, noctis} 'night' (Lewis \& Short): to stay through the night.
\end{cousinweb}

\begin{grammarlens}[title={What you've built}]
One more word for telling, explaining and describing at B1.
\end{grammarlens}

\subsection*{Your turn}
\begin{itemize}
\raggedright
  \item \emph{Say it:} \emph{pernoctō, pernoctāre}
  \item \emph{Say it:} \emph{pernoctō, pernoctāre}, once more
  \item \emph{Say it:} a sentence of your own with \emph{pernoctō, pernoctāre}
  \item \emph{Recall:} say the Latin for a roadside inn, then the Latin for to lodge, then say \emph{pernoctō, pernoctāre} again
\end{itemize}

\begin{tcolorbox}[breakable,colback=teal!4,colframe=teal!35!black,title={Before you move on}]
What does pernoctō, pernoctāre mean? (\textquotedblleft{}To spend the night, to stay all night\textquotedblright{}.) Say it once more.
\end{tcolorbox}
\section[stabulārius, stabulāriī]{stabulārius, stabulāriī — a stable-keeper, an innkeeper (of a stable inn)}

\label{lesson:LA-C445-stabularius}



\begin{quote}
Before the new one: say the Latin for to lodge, then the Latin for to spend the night.
\end{quote}

\subsection*{You'll want to know: stabulārius, stabulāriī}
\textbf{stabulārius, stabulāriī} — \textquotedblleft{}a stable-keeper, an innkeeper (of a stable inn)\textquotedblright{}.

\begin{quote}
Stabulārius equōs viātōrum cūrat.
\end{quote}

\begin{quote}
\emph{The stable-keeper looks after the travellers' horses.}
\end{quote}

Second declension, masculine: \emph{stabulārius, stabulāriī m.} The keeper of a \emph{stabulum}, the cheapest kind of roadside lodging.

\begin{cousinweb}
\emph{stabulārius} is built on \emph{stabulum} 'stable, lodging' with the suffix \emph{-ārius} 'one who works with' (Lewis \& Short).
\end{cousinweb}

\begin{grammarlens}[title={What you've built}]
One more word for telling, explaining and describing at B1.
\end{grammarlens}

\subsection*{Your turn}
\begin{itemize}
\raggedright
  \item \emph{Say it:} \emph{stabulārius, stabulāriī}
  \item \emph{Say it:} \emph{stabulārius, stabulāriī}, once more
  \item \emph{Say it:} a sentence of your own with \emph{stabulārius, stabulāriī}
  \item \emph{Recall:} say the Latin for to lodge, then the Latin for to spend the night, then say \emph{stabulārius, stabulāriī} again
\end{itemize}

\begin{tcolorbox}[breakable,colback=teal!4,colframe=teal!35!black,title={Before you move on}]
What does stabulārius, stabulāriī mean? (\textquotedblleft{}A stable-keeper, an innkeeper (of a stable inn)\textquotedblright{}.) Say it once more.
\end{tcolorbox}
\section[hospitālis, hospitāle]{hospitālis, hospitāle — hospitable; of a guest or host}

\label{lesson:LA-C445-hospitalis}



\begin{quote}
Before the new one: say the Latin for to spend the night, then the Latin for a stable-keeper.
\end{quote}

\subsection*{You'll want to know: hospitālis, hospitāle}
\textbf{hospitālis, hospitāle} — \textquotedblleft{}hospitable; of a guest or host\textquotedblright{}.

\begin{quote}
Caupō hospitālis est et viātōribus bonam cēnam dat.
\end{quote}

\begin{quote}
\emph{The innkeeper is hospitable and gives the travellers a good dinner.}
\end{quote}

An adjective of two endings, from \emph{hospes} 'host, guest'. English \emph{hospital} and \emph{hospitable} come from the same family.

\begin{cousinweb}
\emph{hospitālis} is built on \emph{hospes, hospitis} 'host, guest', from PIE \emph{g\textsuperscript{h}óspot-} 'master of guests' (EtymDB; Lewis \& Short). English \emph{hospital} comes from it through French.
\end{cousinweb}

\begin{grammarlens}[title={What you've built}]
One more word for telling, explaining and describing at B1.
\end{grammarlens}

\subsection*{Your turn}
\begin{itemize}
\raggedright
  \item \emph{Say it:} \emph{hospitālis, hospitāle}
  \item \emph{Say it:} \emph{hospitālis, hospitāle}, once more
  \item \emph{Say it:} a sentence of your own with \emph{hospitālis, hospitāle}
  \item \emph{Recall:} say the Latin for to spend the night, then the Latin for a stable-keeper, then say \emph{hospitālis, hospitāle} again
\end{itemize}

\begin{tcolorbox}[breakable,colback=teal!4,colframe=teal!35!black,title={Before you move on}]
What does hospitālis, hospitāle mean? (\textquotedblleft{}Hospitable; of a guest or host\textquotedblright{}.) Say it once more.
\end{tcolorbox}
